\documentclass[11pt]{article}
\usepackage[margin=1in]{geometry}
\usepackage{amsmath,amssymb,amsthm}
\usepackage{graphicx}
\usepackage[hidelinks]{hyperref}

\newtheorem{theorem}{Theorem}
\newtheorem{lemma}{Lemma}
\newtheorem{proposition}{Proposition}

\newcommand{\E}{\mathbb E}
\newcommand{\R}{\mathbb R}
\newcommand{\K}{\mathsf K}

\title{Exact Product-Measure Constants in a Question on Ramsey-Type Spheres}
\author{Run fa\_banach\_001, agent\_lane\_13}
\date{2026-06-14}

\begin{document}
\maketitle

\section*{Status}

This is a partial solution packet for arXiv:2308.07668, ``On Ramsey-type
properties of the distance in nonseparable spheres'' by Piotr Koszmider.  It
settles the exact values of the constants \(c_n\) and \(c\) appearing in
Proposition 55 of the source paper.  The endpoint dichotomy question for
\(S_{L_p(\{0,1\}^{\kappa})}\) remains separate.

The last finite case, \(1<p<2\), is obtained by identifying the relevant finite
atom inequality with the largest Rayleigh quotient of the complete graph
\(p\)-Laplacian.  The complete graph spectrum is due to Amghibech, and is used
here in the form quoted in Zhang's Section 4.3 \cite{Zhang2023}.

\section*{Source Evidence}

The following crop is from page 25 of the source PDF.

\begin{center}
\includegraphics[width=\textwidth]{figures/open_problem_crop.png}
\end{center}

The relevant source sentence is:
\[
\text{``We do not know the exact value of }c_n\text{ or of }c\text{ in the above proof or if }S_{L_p}\text{ is dichotomous.''}
\]

\section*{Source Notation}

Let \(\mu\) be the product probability measure on \(\{0,1\}^{\kappa}\), where
\(\kappa\leq2^\omega\) is uncountable.  In Proposition 55 of the source paper,
for \(A\in[\kappa]^n\), the map \(T_A\) sends a function
\(F\in L_p(\{0,1\}^n)\) to the function on \(\{0,1\}^{\kappa}\) depending on
the coordinates in \(A\).  The constant \(c_n\) is the supremum of
\[
 \|T_A(F)-T_B(F)\|_{L_p}
\]
over all unit vectors \(F\in L_p(\{0,1\}^n)\) and all admissible pairs
\(A,B\in[\kappa]^n\) whose common coordinates occur in the same positions.
Finally \(c=\sup_n c_n\).

\section*{Proof Intuition}

After the common coordinates of \(A\) and \(B\) are fixed, the two random
variables \(T_A(F)\) and \(T_B(F)\) are independent copies of the same smaller
finite random variable.  Thus the problem is controlled by finite complete
graph quotients
\[
 \sup_{x\neq0}\frac{\sum_{i,j}|x_i-x_j|^p}{\sum_i |x_i|^p}.
\]
For \(1<p<2\), the largest quotient is not the limiting negative-type value:
the finite atom mass forces a strict gap.  Amghibech's spectrum for the
complete graph \(p\)-Laplacian gives exactly this gap, and the largest
eigenvalue is the one corresponding to one positive atom, one negative atom,
and all other atoms zero.  For \(p>2\), the same spectrum is maximized by a
balanced positive/negative split; because \(2^n\) is even, this gives the
Rademacher value at every finite level.

\begin{lemma}[complete graph quotient]\label{lem:complete}
Let \(N\geq2\), and for \(1<p<\infty\) define
\[
 M_{N,p}:=\sup_{x\in\R^N\setminus\{0\}}
 \frac{\sum_{i,j=1}^N |x_i-x_j|^p}{\sum_{i=1}^N |x_i|^p}.
\]
Then
\[
M_{N,p}=2N-4+2^p,\qquad 1<p\leq2.
\]
If \(2\leq p<\infty\) and \(N\) is even, then
\[
M_{N,p}=N2^{p-1}.
\]
\end{lemma}

\begin{proof}
For \(p=2\), the identity
\[
\sum_{i,j}|x_i-x_j|^2=2N\sum_i|x_i-\bar x|^2\leq 2N\sum_i|x_i|^2
\]
gives \(M_{N,2}=2N\), and both displayed formulas agree with this value.

Assume first that \(p\neq2\).  Let \(K_N\) be the complete graph on \(N\)
vertices.  Zhang quotes Amghibech's Theorem 6 as follows: for \(N\geq3\), the
nonzero normalized \(p\)-Laplacian eigenvalues of \(K_N\) are
\[
 \lambda_{i,j}=
 \frac{1}{N-1}\left(
 N-i-j+\left(i^{1/(p-1)}+j^{1/(p-1)}\right)^{p-1}
 \right),
 \qquad i,j\in\mathbb Z_+,\quad i+j\leq N.
\]
Here the normalized Rayleigh quotient is
\[
 F_p(x)=
 \frac{\sum_{\{i,j\}\in E(K_N)}|x_i-x_j|^p}
 {(N-1)\sum_i |x_i|^p}.
\]
Since \(F_p\) attains its maximum on the compact \(\ell_p\)-sphere and a
maximizer is a critical point, \(\sup F_p\) is the largest value among the
\(\lambda_{i,j}\).  Also
\[
 M_{N,p}=2(N-1)\sup F_p,
\]
because \(M_{N,p}\) uses ordered pairs while \(F_p\) uses unordered edges.

Put \(q=1/(p-1)\).  If \(1<p<2\), then \(q>1\).  The expression to maximize is
\[
 H(i,j)=N-i-j+\|(i,j)\|_q.
\]
For \(i,j\geq1\), the quantity \(i+j-\|(i,j)\|_q\) is strictly increasing in
each variable, since
\[
 \frac{\partial}{\partial i}\left(i+j-(i^q+j^q)^{1/q}\right)
 =1-\left(\frac{i^q}{i^q+j^q}\right)^{(q-1)/q}>0.
\]
Thus the maximum is attained at \(i=j=1\), and
\[
\sup F_p=\frac{N-2+2^{p-1}}{N-1}.
\]
Hence \(M_{N,p}=2N-4+2^p\).

If \(p>2\), then \(0<q<1\).  For a fixed \(s=i+j\), concavity of \(t^q\)
implies that \((i^q+j^q)^{1/q}\) is maximized when \(i\) and \(j\) are as
balanced as possible.  When \(N\) is even, taking \(s=N\) and \(i=j=N/2\)
gives
\[
 H(i,j)=\left(2(N/2)^q\right)^{1/q}=N2^{p-2}.
\]
For any \(s\leq N\), the same balanced estimate gives
\[
 H(i,j)\leq N-s+2^{p-2}s,
\]
which is maximized at \(s=N\) because \(2^{p-2}>1\).  Therefore
\(\sup F_p=N2^{p-2}/(N-1)\), and \(M_{N,p}=N2^{p-1}\).

For \(N=2\), the same values follow directly from the two-point quotient, so
the proof is complete.
\end{proof}

\begin{lemma}[conditioning reduction]\label{lem:conditioning}
Let \(N=2^n\).  For the source constants \(c_n\),
\[
 c_n^p\leq \frac{M_{N,p}}{N}.
\]
\end{lemma}

\begin{proof}
Fix \(F\in L_p(\{0,1\}^n)\) and an admissible pair \(A,B\).  Suppose that
\(r=|A\cap B|\), with common coordinates in the same positions, and condition
on these \(r\) common coordinates.  Given the common coordinates,
\(T_A(F)\) and \(T_B(F)\) are iid copies of the same function on a cube with
\(2^{n-r}\) atoms.  Lemma \ref{lem:complete} applied to this smaller finite
cube gives
\[
 \E\left(|T_A(F)-T_B(F)|^p\mid A\cap B\right)
 \leq d_{n-r,p}\,
 \E\left(|T_A(F)|^p\mid A\cap B\right),
\]
where
\[
d_{m,p}=
\begin{cases}
2-(4-2^p)/2^m, & 1<p\leq2,\ m\geq1,\\
2^{p-1}, & 2\leq p<\infty,\ m\geq1,\\
0, & m=0.
\end{cases}
\]
For \(1<p\leq2\), the numbers \(d_{m,p}\) increase with \(m\); for
\(p\geq2\), they are constant for \(m\geq1\).  Therefore
\(d_{n-r,p}\leq M_{2^n,p}/2^n\).  Averaging over the common coordinates and
normalizing \(\|F\|_p=1\) proves the claim.
\end{proof}

\begin{theorem}[exact constants]\label{thm:constants}
For Proposition 55 of the source paper, for every \(n\geq1\),
\[
 c_n^p =
 \begin{cases}
 2-\dfrac{4-2^p}{2^n}, & 1<p\leq2,\\[6pt]
 2^{p-1}, & 2\leq p<\infty.
 \end{cases}
\]
Consequently,
\[
 c=\sup_n c_n =
 \begin{cases}
 2^{1/p}, & 1<p\leq2,\\
 2^{1-1/p}, & 2\leq p<\infty.
 \end{cases}
\]
\end{theorem}

\begin{proof}
The upper bounds follow from Lemmas \ref{lem:complete} and
\ref{lem:conditioning}, with \(N=2^n\).

For the lower bounds, choose disjoint \(A,B\).  If \(1<p\leq2\), let \(F\) take
the value \(a=(2^{n-1})^{1/p}\) on one atom of \(\{0,1\}^n\), the value \(-a\)
on another atom, and \(0\) on all remaining atoms.  Then \(\|F\|_p=1\), and
for independent copies \(U,V\) of \(F\),
\[
\E|U-V|^p
 =\frac{2(2a)^p+4(2^n-2)a^p}{2^{2n}}
 =2-\frac{4-2^p}{2^n}.
\]
Thus \(c_n^p\) is at least the stated value.

If \(p\geq2\), let \(F\) be a Rademacher function depending on one coordinate.
Then independent copies \(U,V\) satisfy \(\E|U-V|^p=2^{p-1}\).  Hence
\(c_n^p\geq2^{p-1}\), and the upper bound gives equality.

Taking the supremum over \(n\) gives the displayed value of \(c\).
\end{proof}

\section*{Verification Notes}

The auxiliary script \texttt{code/check\_finite\_atom\_conjecture.py} was used
before the graph-spectrum identification to test the finite \(1<p<2\) formula
for small \(N\).  It is not part of the proof.  The proof now rests on the
conditioning reduction and on Amghibech's complete-graph \(p\)-Laplacian
spectrum as quoted in \cite[Section 4.3]{Zhang2023}.

\section*{Scope and Limitations}

The exact constants \(c_n\) and \(c\) are settled.  The remaining source
question is whether the unit sphere \(S_{L_p(\{0,1\}^{\kappa})}\) is
dichotomous.  For \(p=2\), the earlier Hilbertian packet records the negative
answer via the Erdos--Preiss and Preiss--Rodl theorem.  For \(p\neq2\), the
endpoint dichotomy problem is not resolved here.

\section*{Literature and Novelty Notes}

The finite \(1<p<2\) formula was first found here numerically, but the proof
uses an existing graph \(p\)-Laplacian theorem of Amghibech.  A web search for
``complete graph p-Laplacian eigenvalues'', ``largest eigenvalue graph
p-Laplacian complete graph'', and exact pairwise-sum phrases led to Zhang's
arXiv:2110.06054 paper, whose Section 4.3 quotes the decisive Amghibech
formula.  I did not locate a freely downloadable copy of Amghibech's original
2003 paper during the bounded search, so Zhang's arXiv PDF is included as the
supporting source for the quoted formula.

\section*{Human Review Recommendation}

Reviewers should first check the factor \(M_{N,p}=2(N-1)\sup F_p\), since this
is the only normalization-sensitive step.  Then check that the integer
maximization of Amghibech's formula is correct in the two regimes \(1<p<2\)
and \(p>2\), and that the conditioning argument covers every admissible
overlap pattern from Proposition 55.

\begin{thebibliography}{9}

\bibitem{Koszmider2024}
P.~Koszmider,
\emph{On Ramsey-type properties of the distance in nonseparable spheres},
arXiv:2308.07668.

\bibitem{Amghibech2003}
S.~Amghibech,
\emph{Eigenvalues of the discrete \(p\)-Laplacian for graphs},
Ars Combinatoria 67 (2003), 283--302.

\bibitem{Zhang2023}
D.~Zhang,
\emph{Homological eigenvalues of graph \(p\)-Laplacians},
Journal of Topology and Analysis (2023), arXiv:2110.06054.

\bibitem{Schoenberg1938}
I.~J. Schoenberg,
\emph{Metric spaces and positive definite functions},
Transactions of the American Mathematical Society 44 (1938), 522--536.

\end{thebibliography}

\end{document}
